\section{Online Implications}\label{sec:online_implications}

The horizon law identifies the quantities that any adaptive memory rule must
track. A policy that updates memory online needs a finite-sample regime,
represented by the exponent $a$ or by an operational exponent
$a_{\varepsilon}$, together with a local roughness description of the recent
path through the pair $(H,\zeta)$. Once these objects are available, the
structural scale $n_\star$ and the useful-memory region $\mathcal U_\delta$
follow directly from the horizon law.

In operational geometries, local support structure enters before roughness
estimation. The embedded fixed-$\varepsilon$ theory already shows that intrinsic
complexity controls which finite-sample regime is available through the
threshold $\alpha>k/2$. A local proxy such as TwoNN therefore serves as a
regime-selection statistic: it does not estimate the horizon by itself, but it
selects the finite-sample law to which the horizon formula should be attached.

The roughness side is a separate statistical task. In deterministic H\"older
paths, $H$ and $\zeta$ belong to an envelope governing how quickly staleness can
accumulate with lag. A direct log--log fit to one short empirical sequence of
transport discrepancies is generally unstable, because single-path local motion,
finite-sample noise, and short-lag discreteness all distort the envelope. The
relevant online object is therefore not a raw local slope, but a stabilized
estimate built from block averages, activity proxies, or other aggregated recent
discrepancies.

This separation gives the minimal architecture of a horizon-informed adaptive
policy. First estimate the local regime that determines the finite-sample term.
Then estimate a recent roughness envelope for the staleness term. Finally map
these estimates to $n_\star$ or to a tolerance band $\mathcal U_\delta$, and use
that region as the admissible memory range. Hysteresis, band projection, or
validation scoring can stabilize the resulting controller, but those mechanisms
operate downstream of the horizon law itself.

The online problem is therefore structurally well posed. The horizon law
supplies the target for adaptation; the remaining task is to estimate its
ingredients with guarantees strong enough to keep a data-driven memory rule
inside the useful-memory region.
